A class is built on another by naming it.
The statement \codehl{class D(B)} defines a class $D$
whose linearisation is $D$ followed by the linearisation of $B$;
$D$ is a \emph{subclass} of $B$, and $B$ is the \emph{base class} of $D$.
Every class of the package names at most one base class, and we assume so throughout.
By Definition \ref{def.attributeReference},
a name that $D$ does not bind is found in $B$, and $D$ is said to \emph{inherit} it;
a name that $D$ binds is found before the one of $B$, and $D$ is said to \emph{override} it.
An instance $x$ of $D$ counts as one of $B$:
\codehl{isinstance(x, B)} is true whenever $B$ is in the linearisation of the class of $x$.

\begin{prop}\label{prop.lateBinding}
 Let \codehl{f} be a function bound in the class $B$,
 and $x$ an instance of a subclass $D$ of $B$ that inherits it.
 In the call \codehl{x.f(u)}, a reference \codehl{self.g} in the body of the function
 is the reference \codehl{x.g}:
 it is resolved along the linearisation of $D$, and not along that of $B$.
\end{prop}
\begin{proof}
 By Definition \ref{def.attributeReference} the call passes $x$ as first argument,
 so \codehl{self} is a name of $x$, and $x$ is an instance of $D$.
\end{proof}

So a method that $D$ inherits calls the functions that $D$ overrides.
A subclass changes a part of what its base class does
by overriding the functions in which that part is written;
it inherits the rest, which then runs on what the subclass changed.
This is the construction of Section \ref{sec.fasterBook}.
The matching of Listing \ref{listing.submit} is written in \codehl{AggregateBook}
against \codehl{self.best\_price} and \codehl{self.set\_size},
and \codehl{CachedBestBook}, in Listing \ref{listing.cachedBestBook},
overrides those two and binds nothing else of the fold.

Inside a function of the class $C_k$ of a linearisation,
\codehl{super().name} is the reference \codehl{self.name}
with the search of part \emph{i.} begun at $C_{k+1}$.
It is how a function calls the one it overrides:
in \codehl{CachedBestBook}, the call \codehl{super().set\_size(direction, price, size)}
is \codehl{AggregateBook.set\_size(self, direction, price, size)}.
The constructor of \codehl{CachedBestBook} calls that of \codehl{AggregateBook},
which binds the two dictionaries, and adds the cache.
Its \codehl{set\_size} writes the dictionary through the function it overrides,
and then brings the cache up to date.
Its \codehl{best\_price} returns the cache, and does not call the scan it overrides.

\begin{lstlisting}[caption={\protect\codehl{unito26.lob.orderbook.CachedBestBook}, abridged: the constructor and the two overrides.}, label=listing.cachedBestBook]
class CachedBestBook(AggregateBook):
    # ...
    def __init__(self, strict: bool = False):
        super().__init__(strict)
        self._cached: dict[int, int | None] = {BUY: None, SELL: None}

    def set_size(self, direction: int, price: int, size: int) -> None:
        super().set_size(direction, price, size)
        self._cached[direction] = self._best_after(direction, price, size)
    # ...
    def best_price(self, direction: int) -> int | None:
        return self._cached[direction]
\end{lstlisting}

\begin{remark}\label{remark.constructorsBindTheState}
 A subclass inherits the attributes of its base class.
 The attributes of an instance are not inherited:
 they are bound by the constructors that run,
 and a constructor that does not call \codehl{super().\_\_init\_\_}
 leaves the instance without those of the base class.
 It is in this sense that \codehl{TickArrayBook} deletes the two dictionaries it inherits,
 in Listing \ref{listing.tickArrayInit}:
 the constructor of \codehl{AggregateBook} has just bound them.
\end{remark}

Listing \ref{listing.lateBinding} writes a subclass.
\codehl{CountingBook} overrides the constructor, to add a count,
and \codehl{best\_price}, to increase the count and then call the function it overrides.
It binds neither \codehl{apply} nor \codehl{submit}.
The market buy of the worked example of Section \ref{sec.fasterBook} is applied to it:
the inherited matching trades at two prices,
and by Proposition \ref{prop.lateBinding} it asks for the best ask through the override,
once for each.

\begin{lstlisting}[caption={An inherited method calls the function that the subclass overrides.}, label=listing.lateBinding]
from unito26.lob.messages import BUY, SELL, market_order
from unito26.lob.orderbook import AggregateBook


class CountingBook(AggregateBook):
    def __init__(self, strict):
        super().__init__(strict)
        self.lookups = 0

    def best_price(self, direction):
        self.lookups += 1
        return super().best_price(direction)


assert CountingBook.__mro__ == (CountingBook, AggregateBook, object)
assert CountingBook.submit is AggregateBook.submit

book = CountingBook(strict=False)
for price, size in ((101, 100), (102, 80), (105, 20)):
    book.set_size(SELL, price, size)
book.apply(market_order(1.0, 150, BUY), record=False)
assert book.lookups == 2
assert book.size_at(SELL, 102) == 30
\end{lstlisting}

Inheritance gives $D$ the attributes of $B$,
and it guarantees nothing about what the overrides of $D$ do.

\begin{defi}[{\citealp{LW94beh}}]\label{def.substitution}
 A subclass $D$ \emph{substitutes} for its base class $B$
 if every statement that holds of the instances of $B$,
 and is written in the attributes of $B$,
 holds of the instances of $D$.
\end{defi}

For the books the statement is the fold \eqref{eq.bookFold}.
The storage of a book stands for a configuration of Section \ref{sec.orderDrivenMarkets}.
Write $\configurationOf(s)$ for the configuration that a storage $s$ stands for,
which \codehl{levels\_map} returns one side at a time,
and $F_D(s, m)$ for the storage that \codehl{apply} leaves
when a book of the class $D$, in the storage $s$, is given the message $m$.
Then $D$ substitutes for \codehl{AggregateBook} in the fold when
\begin{equation}\label{eq.commutingFold}
 \configurationOf\big(F_D(s, m)\big) = F\big(\configurationOf(s), m\big)
\end{equation}
for every storage $s$ that satisfies the invariants of $D$ and every message $m$,
where $F$ is the transition of \eqref{eq.bookFold}.
The map $\configurationOf$ and the condition \eqref{eq.commutingFold}
are those of \cite{Hoa72pro};
for \codehl{AggregateBook} the storage is the configuration,
and $\configurationOf$ is the identity.
Equation \eqref{eq.commutingFold} is the statement of Section \ref{sec.fasterBook}
that two books which share their matching pass through the same configurations,
and the test module named there checks it after every message.

\begin{remark}\label{remark.limitsOfSubstitution}
 Three things lie outside \eqref{eq.commutingFold}.
 The cost of an operation is not a statement about configurations,
 and Table \ref{tab.fasterBookCosts} is the reason for having more than one class.
 \codehl{TickArrayBook} refuses a write outside its band,
 and \codehl{BitmapBook} one below its floor,
 so each satisfies \eqref{eq.commutingFold} for the messages whose prices it can hold.
 And the constructors differ:
 \codehl{TickArrayBook} takes a floor and a width,
 where \codehl{AggregateBook} takes neither.
 Code that builds a book of a class it is handed
 therefore calls the class method \codehl{for\_prices},
 which these two classes override to compute their own arguments,
 or \codehl{from\_levels}, which calls it.
\end{remark}

A class is an object, by Remark \ref{remark.completeRule},
and \codehl{AXIS\_B\_VARIANTS} in \codehl{unito26.lob.orderbook}
is the tuple of the five book classes of Section \ref{sec.fasterBook}.
Listing \ref{listing.fiveBooks} builds a book of each
through the one function \codehl{from\_levels} of Listing \ref{listing.fromLevels},
applies the same market order to each,
and reads the same configuration off each.
It then checks, on the attributes of the four subclasses,
that none binds the matching and that each binds the write and the lookup.

\begin{lstlisting}[caption={Five classes, one matching, one configuration.}, label=listing.fiveBooks]
from unito26.lob.messages import BUY, SELL, market_order
from unito26.lob.orderbook import AXIS_B_VARIANTS, AggregateBook

asks = {101: 100, 102: 80, 105: 20}
for cls in AXIS_B_VARIANTS:
    book = cls.from_levels({99: 100}, asks)
    assert type(book) is cls and isinstance(book, AggregateBook)
    book.apply(market_order(1.0, 150, BUY), record=False)
    assert book.levels_map(SELL) == {102: 30, 105: 20}

matching = {"apply", "submit", "withdraw"}
for cls in AXIS_B_VARIANTS[1:]:
    assert not matching & set(vars(cls))
    assert {"set_size", "best_price"} <= set(vars(cls))
\end{lstlisting}
